% flavour text, presentation of the problem, may contain images and quotes
% image:
%\begin{figure}[h]
%\center\includegraphics[width=0.3\textwidth]{name.png}
%\end{figure}

\problemname{Ratatöskr}

Ratatöskr is a squirrel that lives in a giant (but finite) mythical tree called Yggdrasil. He likes to gossip, which sets the other inhabitants of the tree against each other. Ratatöskr is thus hunted by the two ravens of Odin, which are called Hugin and Munin, to bring him to justice.

The tree is made up of nodes connected by branches. Initially, the ravens and the squirrel sit on three different nodes. % [Alternatively: the ravens are in the air in the beginning?]
Now the following happens repeatedly:
\begin{itemize}
\item On Odin's signal, one of the ravens launches into the air and flies to another node of the tree (or possibly back to its previous position), while the other stays where it is. %[Alternatively: both can move?] 
\item During this maneuver, Ratatöskr can travel along the branches to reach another node, but may not pass through a node where a raven sits. He is much quicker than the ravens and will reach his destination before the flying raven lands.
\end{itemize}
Ratatöskr gets captured if one of the ravens flies to his position and there is no other node he can escape to.
%We may safely assume that he will not pick a destination where a raven chooses to land, if this is possible. However, if Ratatöskr cannot move at all due to a sitting raven, he gets caught if the other raven chooses to fly to his position.

Help Odin determine an optimal strategy for capture, i.e. the minimum number of signals he has to give until Ratatöskr is guaranteed to be captured by a raven.
%Output \texttt{impossible} if Ratatöskr can escape them indefinitely.
%The two ravens can be assumed to never share a node, as this is never beneficial, and there are at least three nodes (so it is possible).

\section*{Input}
The input consists of:
\begin{itemize}
	%\item one line with two integers $n$ ($1\leq n\leq 10^3$) and $m$ ($1\leq m\leq 5\cdot 10^3$), where $n$ is the amount of \ldots
	\item one line with a single integer $n$ ($3\leq n\leq 80$), the number of nodes in the tree. The nodes are labeled $1, \dots, n$.
	\item one line with a single integer $r$ ($1 \leq r \leq n$), the initial position of the squirrel Ratatöskr.
	\item one line with a single integer $h$ ($1 \leq h \leq n$), the initial position of the raven Hugin.
	\item one line with a single integer $m$ ($1 \leq m \leq n$), the initial position of the raven  Munin.
	\item $n-1$ lines each containing two integers $i$ and $j$ ($1 \leq i < j \leq n$), indicating a branch between nodes $i$ and $j$.
\end{itemize}

The positions $r$, $h$ and $m$ are distinct. There is at most one branch between any two nodes and every node is reachable from every other node by a sequence of branches.


%The input consists of:
%\begin{itemize}
%%\item one line with two integers $n$ ($1\leq n\leq 10^3$) and $m$ ($1\leq m\leq 5\cdot 10^3$), where $n$ is the amount of \ldots
%\item one line with an integer $3\leq n\leq 80$, the number of nodes in the tree. The nodes are labeled $0, \dots, n-1$.
%\item one line with an integer $0 \leq r \leq n-1$, the initial position of the squirrel Ratatöskr.
%\item one line with an integer $0 \leq h \leq n-1$, $h \neq r$, the initial position of the raven Hugin.
%\item one line with an integer $0 \leq m \leq n-1$, $h \neq m \neq r$, the initial position of the raven  Munin.
%\item $n-1$ lines, each containing two integers $0 \leq i < j \leq n-1$, indicating a connection between two nodes $i$ and $j$. (There is at most one connection between two nodes.)
%\end{itemize}

\section*{Output}
% detailed description of the output format
One line containing an integer $s$, the number of signals that the ravens need to capture Ratatöskr in an optimal strategy. If Ratatöskr can escape them indefinitely, output \texttt{impossible}.
%For each test case, print a line containing ``Case \#$i$: Hello $s$!'' where $i$ is its number, starting at 1. Each line of the output should end with a line break.
